% Варіант підрозділу 4.1 для обговорення; головний документ не змінено.
% Фрагмент використовує позначення й оточення main-SMPT_UK.tex.

\subsection{Точність для мережі з рівномірними міжспільнотними зв’язками}
\label{sec:surrogate}

Тепер покажемо, яку мережу точно моделює ідеальний процес,
описаний у Розділі~\ref{sec:method}. Для цього побудуємо
сурогатну мережу з рівномірними міжспільнотними зв’язками
$\tilde G$. Внутрішні контакти та їхні інтенсивності залишимо
такими самими, як у початковій мережі $G$. Між кожними двома
зв’язаними спільнотами замість справжніх каналів $E_{ij}$
використаємо всі канали $V_i\times V_j$: кожен вузол джерела
буде з’єднаний із кожним вузлом одержувача.

Відома загальна вага зв’язку $W_{ij}$ не змінюється:
розподіляємо її порівну між усіма $N_iN_j$ парами вузлів,
де $N_i$ та $N_j$ --- кількості вузлів у відповідних спільнотах.
Тому кожен новий канал має інтенсивність
$\tilde\beta_{uv}=\beta_{\mathrm{out}}W_{ij}/(N_iN_j)$,
а сума інтенсивностей каналів дорівнює
\[
\sum_{(u,v)\in V_i\times V_j}\tilde\beta_{uv}
=\beta_{\mathrm{out}}W_{ij}.
\]
Таким чином, зберігаємо загальну вагу зв’язку, доступну
агрегованому інтерфейсу, але змінюємо конкретні вузли,
між якими можливе передавання інфекції. Наступне твердження
пов’язує цю мережу з правилами двомасштабного методу.

\begin{proposition}[Точність для мережі з рівномірними міжспільнотними зв’язками]
\label{prop:surrogate}
За рівномірного вибору цілі імпорту серед сприйнятливих вузлів
ідеальний двомасштабний процес має той самий розподіл,
що й повномережевий Gillespie на побудованій вище мережі
$\tilde G$, якщо обидва процеси починаються з однакового стану.
\end{proposition}

\begin{proof}
Локальні зараження та видалення мають однакові інтенсивності,
оскільки внутрішні контакти не змінено. Залишається порівняти
зараження з інших спільнот.

Для сприйнятливого вузла $v\in V_j$ імовірність вибору ціллю
імпорту дорівнює $1/s_j(x)$. Множимо інтенсивність імпортів
із кожної спільноти $i$ на цю імовірність і підсумовуємо:
\[
\begin{aligned}
\sum_{i\neq j}\lambda_{ij}(x)\frac{1}{s_j(x)}
&=\sum_{i\neq j}\beta_{\mathrm{out}}W_{ij}
\frac{y_i(x)}{N_i}\frac{s_j(x)}{N_j}\frac{1}{s_j(x)}\\
&=\sum_{i\neq j}\sum_{u\in V_i}
\frac{\beta_{\mathrm{out}}W_{ij}}{N_iN_j}\ind\{x_u=I\}.
\end{aligned}
\]
Останній вираз --- сума інтенсивностей усіх каналів $\tilde G$,
які ведуть від інфікованих вузлів інших спільнот до $v$.
Отже, інтенсивність зараження $v$ однакова в обох процесах.
Якщо $s_j(x)=0$, сприйнятливих цілей немає й імпорт неможливий.
Оскільки всі інтенсивності переходів збігаються, за однакового
початкового стану збігаються й розподіли процесів. \qed
\end{proof}

Якщо міжспільнотні контакти в $G$ уже з’єднують усі пари
вузлів відповідних спільнот і мають однакову інтенсивність
для кожної пари спільнот, то $\tilde G=G$. У цьому окремому
випадку ідеальний двомасштабний процес точний і для початкової
мережі. Внутрішні мережі при цьому не мусять бути повними.

У загальному випадку $G$ і $\tilde G$ відрізняються
міжспільнотними контактами. Однакова сума інтенсивностей
каналів не означає однакової поточної інтенсивності зараження:
її визначають лише контакти від інфікованих вузлів до
сприйнятливих. Тому заміна контактів може змінити результати
навіть за точного визначення часу подій. Саме ці відмінності,
а також додаткову похибку часової сітки в реалізації,
розглядаємо далі.
